\section*{Problem 1}

给定二次型
\[
Q \left(x_1, x_2, x_3\right) = 2x_1^2 + 3x_2^2 + x_3^2 + 4x_1x_2 - 2x_1x_3 + 6x_2x_3.
\]
写出其对应的对称矩阵.

\begin{proof}

设 $\boldsymbol{x} = \left(x_1, x_2, x_3\right)^\top$，则有
\[
\begin{aligned}
Q \left(\boldsymbol{x}\right) &= 2 x_1^2 + 3 x_2^2 + x_3^2 + 2 x_1x_2 + 2 x_2x_1 - x_1x_3 - x_3x_1 + 3x_2x_3 + 3x_3x_2 \\
&= \left(x_1, x_2, x_3\right) \begin{bmatrix}
    2 & 2 & -1 \\
    2 & 3 & 3 \\
    -1 & 3 & 1 \\
\end{bmatrix} \left(x_1, x_2, x_3\right)^\top = \boldsymbol{x}^\top A \boldsymbol{x}
\end{aligned}
\]
于是 $\boxed{A = \begin{bmatrix}
    2 & 2 & -1 \\
    2 & 3 & 3 \\
    -1 & 3 & 1 \\
\end{bmatrix}}.$

\end{proof}

\section*{Problem 2}

判断下列二次型的正定性：

\begin{enumerate}[]
    \item $Q_1 \left(x, y\right) = x^2 + 2xy + 2y^2, $
    \item $Q_2 \left(x, y, z\right) = x^2 + y^2 + z^2 - xy - xz - yz.$
\end{enumerate}

\begin{proof}
对 $Q_1 \left(x, y\right) = x^2 + 2xy + 2y^2 = \left(x, y\right) A_1 \left(x, y\right)^\top$，易见
\[
A_1 = \begin{bmatrix}
    1 & 1 \\
    1 & 2 \\
\end{bmatrix}.
\]
为求 $A_1$ 的特征值，解方程
\[
\det \left(\lambda I - A_1\right) = 0 \iff \det \begin{bmatrix}
    \lambda - 1 & -1 \\
    -1 & \lambda - 2 \\
\end{bmatrix} = \lambda^2 - 3 \lambda + 1 = 0.
\]
即有 $\lambda_1 + \lambda_2 = 3, \lambda_1 \lambda_2 = 1 \iff \lambda_1, \lambda_2 > 0$，因此 $\boxed{Q_1 \textrm{ 是正定的}}$.

对 $Q_2 \left(x, y, z\right) = x^2 + y^2 + z^2 - xy - xz - yz = \left(x, y, z\right) A_2 \left(x, y, z\right)^\top$，易见
\[
A_2 = \begin{bmatrix}
    1 & -1/2 & -1/2 \\
    -1/2 & 1 & -1/2 \\
    -1/2 & -1/2 & 1 \\
\end{bmatrix}.
\]
为求 $A_2$ 的特征值，解方程
\[
\begin{aligned}
\det \left(\lambda I - A_2\right) &= 0 \iff \det \begin{bmatrix}
    \lambda - 1 & 1/2 & 1/2 \\
    1/2 & \lambda - 1 & 1/2 \\
    1/2 & 1/2 & \lambda - 1 \\
\end{bmatrix} = 0, \\
&= \left(\lambda - 1\right) \det \begin{bmatrix}
    \lambda - 1 & 1/2 \\
    1/2 & \lambda - 1 \\
\end{bmatrix} - \frac{1}{2} \det \begin{bmatrix}
    1/2 & 1/2 \\
    1/2 & \lambda - 1 \\
\end{bmatrix} + \frac{1}{2} \det \begin{bmatrix}
    1/2 & \lambda - 1 \\
    1/2 & 1/2 \\
\end{bmatrix}, \\
&= \left(\lambda - 1\right) \left[\left(\lambda - 1\right)^2 - \frac{1}{4}\right] - \frac{1}{4} \left(\lambda - \frac{3}{2}\right) + \frac{1}{4} \left(\frac{3}{2} - \lambda\right), \\
&= \lambda \left(\lambda^2 - 3\lambda + \frac{9}{4}\right) = \lambda \left(\lambda - \frac{3}{2}\right)^2 = 0 \implies \lambda_i \ge 0.
\end{aligned}
\]
因此 $\boxed{Q_2 \textrm{ 是半正定的}}.$

\end{proof}

\section*{Problem 3}

设 $f(\boldsymbol{x}) = \boldsymbol{x}^\top A \boldsymbol{x}$，定义双线性形
\[
g\left(\boldsymbol{x}, \boldsymbol{y}\right) = \frac{1}{2} \left[f \left(\boldsymbol{x} + \boldsymbol{y}\right) - f\left(\boldsymbol{x}\right) - f\left(\boldsymbol{y}\right)\right].
\]
若
\[
A =
\begin{bmatrix}
2 & 1 \\
1 & 3
\end{bmatrix},
\]
求 $g \left(\boldsymbol{e}_1, \boldsymbol{e}_2\right)$，其中 $\boldsymbol{e}_1 = \left(1, 0\right)^\top, \boldsymbol{e}_2 = \left(0, 1\right)^\top$.

\begin{proof}
为求 $g \left(\boldsymbol{e}_1, \boldsymbol{e}_2\right) = \frac{1}{2} \left[f \left(\boldsymbol{e}_1 + \boldsymbol{e}_2\right) - f\left(\boldsymbol{e}_1\right) - f\left(\boldsymbol{e}_2\right)\right]$，先求
\[
\begin{cases}
f \left(\boldsymbol{e}_1\right) &= \left(1, 0\right) \begin{bmatrix}
2 & 1 \\
1 & 3
\end{bmatrix} \left(1, 0\right)^\top = 2, \\
f \left(\boldsymbol{e}_2\right) &= \left(0, 1\right) \begin{bmatrix}
2 & 1 \\
1 & 3
\end{bmatrix} \left(0, 1\right)^\top = 3, \\
f \left(\boldsymbol{e}_1 + \boldsymbol{e}_2\right) &= \left(1, 1\right) \begin{bmatrix}
2 & 1 \\
1 & 3
\end{bmatrix} \left(1, 1\right)^\top = 7. \\
\end{cases}
\]
于是
\[
g \left(\boldsymbol{e}_1, \boldsymbol{e}_2\right) = \frac{7 - 2 - 3}{2} = \boxed{1}.
\]
\end{proof}

\section*{Problem 4}

将二次型
\[
Q \left(x, y\right) = 2x^2 + 4xy + 5y^2
\]
化为标准形，并写出所用的可逆变换矩阵 $P$.

\begin{proof}
注意到
\[
\begin{aligned}
    Q \left(x, y\right) = \left(x, y\right) \underbrace{\begin{bmatrix}
        2 & 2 \\
        2 & 5 \\
    \end{bmatrix}}_{A} \left(x, y\right)^\top &= 2x^2 + 4xy + 5y^2, \\
    &= 2 \left(x + y\right)^2 + 3 y^2 = \left(x + y, y\right) \underbrace{\begin{bmatrix}
        2 & 0 \\
        0 & 3 \\
    \end{bmatrix}}_{B} \left(x + y, y\right)^\top.
\end{aligned}
\]
引入矩阵 $P^{-1}$ 有
\[
\left(x + y, y\right)^\top = \underbrace{\begin{bmatrix}
    1 & 1 \\
    0 & 1 \\
\end{bmatrix}}_{P^{-1}} \left(x, y\right)^\top.
\]
求解 $P$ 可得
\[
P^{-1} = \begin{bmatrix}
    1 & 1 \\
    0 & 1 \\
\end{bmatrix}, P^{-1}P = I \implies
\boxed{P = \begin{bmatrix}
    1 & -1 \\
    0 & 1 \\
\end{bmatrix}}.
\]
\end{proof}
